\section{离光锥软函数的欧氏化}\label{sec:yzliu-soft}

离光锥软函数 $S(b_\perp,\mu,Y,Y')$ 是两个订书钉形 Wilson 回路的时序真空期望：入射段沿类时速度 $v=\gamma(1,\beta,\mathbf 0_\perp)$，出射段沿 $v'=\gamma'(1,-\beta',\mathbf 0_\perp)$（$\beta=\tanh Y$，$\beta'=\tanh Y'$），除以矩形回路之积的平方根后取 $t,t'\to\infty$（\cref{def:lamet-soft-offlc}，\cref{fig:yzliu-soft-hqet}(a)）。本节说明它为何可以在欧氏空间计算。


\begin{lemma}[部分指数化引理 \tagA]\label{lem:yzliu-partial-exp}
设 $W(C;x'\to x)$ 为沿任意分段光滑路径 $C$ 的基础表示 Wilson 线（\cref{def:prelim-wilson}），$O(x)$、$O'(x')$ 为端点上的定域算符，$F(C,x,x')\equiv\Tr\langle O(x)W(C;x'\to x)O'(x')\rangle$。则在微扰论每一阶（Feynman 规范，维数正规化）
\begin{equation}\label{eq:yzliu-partial-exp}
F(C,x,x')=\ee^{\Phi(C)}\,\Tr\big\langle O(x)\,W(C;x'\to x)\,O'(x')\big\rangle_{\mathrm{2PI}} ,
\end{equation}
其中 $\ee^{\Phi(C)}$ 是线本身迹归一化的真空期望，$\Phi(C)$ 为其全部 web（胶子端点都在 $C$ 上的不可约图）之和，色因子取最大非阿贝尔的修正色因子（\cref{prop:appB-webs}）；$\langle\cdots\rangle_{\mathrm{2PI}}$ 为算符—规范线二粒子不可约顶点 $w_0$ 之和，其色因子同样按修正规则计算（一般不是最大非阿贝尔的）。对带算符插入的闭合回路，同一结论借助复制技巧成立。（本引理只涉及微扰论中的算符与 Wilson 线关联函数。）
\uses{prop:appB-webs,def:prelim-wilson}
\end{lemma}

\begin{proofsketch}
（1）每个图的色因子分解为 $C(G)=\sum_{d}\prod_{i=0}^{n(d)}\tilde C(w_i)$，$d$ 由一个顶点 $w_0$ 与若干 web $w_{i\ge1}$ 组成，这是 Gatheral 与 Frenkel--Taylor 的分解\cite{Gatheral1983,FrenkelTaylor1984}在有端点算符时的推广；（2）固定 $w_0$ 的推广程函恒等式 $\sum_{G\in\mathcal G_d}\mathcal I(G)=S_d^{-1}\prod_{i=0}^{n(d)}\mathcal I(w_i)$，$S_d=\prod_in_i!$\cite[式~(B.5)]{Liu:2022nnk}；（3）对 $d$ 求和，web 部分组成 $\ee^{\Phi(C)}$；$\ee^{\Phi(C)}=\Tr\langle W(C)\rangle$ 由复制技巧（$\langle W\rangle^n=1+n\ln\langle W\rangle+\order{n^2}$，只有 web 的色因子对 $n$ 线性）得到。完整论证见\cite[第~3.2~节与附录~B]{Liu:2022nnk}。
\end{proofsketch}

\paragraph{地位与缺口。}全阶论证 [A]。闭合回路的 web 指数化是标准的全阶论证（\cref{prop:appB-webs}，[A]）；推广到固定的顶点 $w_0$ 只被陈述为“与标准情形类似”，恒等式在 $w_0$ 固定时的证明没有写出\cite[附录~B]{Liu:2022nnk}——这是最弱的一步。写出这一组合证明即可把引理提升为 [R, $R_{\mathrm{PT}}$]。

%%FIG fig:yzliu-soft-hqet

\paragraph{捏点奇点与超软模式。}准 TMD 订书钉（\cref{def:lamet-quasi-tmd}）的两条反平行长边在单胶子交换中给出 $1/[(k^z+\ii0)(k^z-\ii0)]$，在 $k^z=0$ 处被钳住；按 $z$ 坐标统一排序两边的顶点（“单线记号”：一对反平行规范线合成一条量纲为 $\frac32$ 的有效线）后，交换图化为单条规范线的自能图乘以额外的 $1/(-\ii0)$，后者在有限长度 $L$ 下正比于 $L$\cite[式~(41)--(43)]{Liu:2022nnk}。它的意义是一对色源沿虚时间的传播，需要 LSZ 型的截肢。

\begin{proposition}[捏点奇点、超软模式与 $L\to\infty$ 极限 \tagA]\label{prop:yzliu-pinch-ultrasoft}
在微扰论各阶：(a) 对准 TMD 订书钉 $\mathcal W_z$，$\ee^{\Phi(\mathcal W_z)}=\langle0|\mathcal W_z|0\rangle\sim\ee^{-LE(b_\perp)}$，$E=2\delta m+V(b_\perp)$，而矩形回路 $Z_E(2L,b_\perp,\mu)\sim\ee^{-2LE}$，故 $\lim_{L\to\infty}\langle P|\tilde O_z|P\rangle/\sqrt{Z_E}$ 逐阶存在，全部捏点奇点都在 $\ee^{\Phi(\mathcal W_z)}$ 中；(b) 在高阶，单个图出现正比于 $1/L$ 的幂次红外发散（超软模式），这是静态势中 Appelquist--Dine--Muzinich 效应的类比——静态势的伪红外发散自 $\order{\alpha_s^4}$ 起出现\cite{Appelquist:1977es,Brambilla:1999qa}；把偶极子的二粒子不可约自相互作用重求和为有效单线的修饰传播子后，这些发散相消，相关能标是 $V_8-V$，对快度为 $Y$ 的软函数为 $\alpha_sC_A\ee^{-Y}/(2b_\perp)$；依赖 $Y$ 的“超—超软”模式被论证为由色中性与 Ward 恒等式相消；(c) 类时、双重时序的软函数没有捏点奇点且红外有限。（本命题只涉及部分子态与真空中的 Wilson 线关联函数。）
% errtype: E-S
\uses{lem:yzliu-partial-exp,def:lamet-quasi-tmd,def:lamet-soft-offlc}
\end{proposition}

\begin{proofsketch}
（a）由\cref{lem:yzliu-partial-exp}，$\langle P|\tilde O_z|P\rangle=\ee^{\Phi(\mathcal W_z)}\langle P|\tilde O_z|P\rangle_{\mathrm{2PI}}$；单线记号下钳夹只来自偶极子自相互作用的插入，它们全部在 $\ee^{\Phi}$ 中，而类空订书钉相当于一对沿虚时间传播的重夸克\cite[第~3.2~节]{Liu:2022nnk}。（b）与 Lamb 位移类似：逐阶展开修饰传播子 $1/[k\cdot v'-\Sigma(k\cdot v',b_\perp)+\ii0]$ 得到幂次发散，重求和后 $\Sigma(0,b_\perp)\sim\ee^{-Y}\alpha_sC_A/(2b_\perp)$ 起截断作用\cite[第~3.3、4.3~节，式~(109)--(111)]{Liu:2022nnk}。（c）对“4PI” web 的数目归纳\cite[附录~D]{Liu:2022nnk}。
\end{proofsketch}

\paragraph{地位与缺口。}全阶论证 [A]；(a) 是\cref{asm:lamet-staple}(ST2) 的全阶论证。(b) 中重求和后的全阶相消与“超—超软”模式的相消是论证而非证明——最弱的一步（幂次计数层次，误差类型 E-S）；(c) 的归纳假定低阶已无幂次红外发散。原文对超软效应在准 TMD 中开始出现的圈数说法不一，本书只引用静态势的陈述。用矩形回路除去离光锥订书钉自相互作用的做法可追溯到文献\cite{Collins:2008ht,Ji:2018hvs}，(a) 为之提供了全阶依据。升级途径见\cref{rem:yzliu-open}第~(8)~项。

\paragraph{作为形状因子的软函数。}沿速度 $v$ 运动的静态色源的拉氏量 $\psi_v^\dagger\,\ii v\cdot D\,\psi_v$ 给出传播子 $\ii/(v\cdot k+\ii0)$ 与顶点 $\ii gt^av^\mu$，恰是程函规则（\cref{prop:prelim-wilson-rules}）：静态色源的传播子就是一条 Wilson 线。

\begin{theorem}[离光锥软函数是快速色源对的形状因子 \tagA]\label{thm:yzliu-soft-hqet}
设 $\psi_v$、$\eta_v$ 为沿 $v$ 运动的静态色源（基础与反基础表示），$\mathcal L=\psi_v^\dagger\,\ii v\cdot D\,\psi_v+\eta_v^\dagger\,\ii v\cdot D\,\eta_v$（$D$ 取本书的符号），
\begin{equation}\label{eq:yzliu-hqet-J}
J(\vec b_\perp,v,v')=\eta^\dagger_{v'}(\vec b_\perp)\,\eta_v(\vec b_\perp)\,\psi^\dagger_{v'}(0)\,\psi_v(0)
\end{equation}
为 $t=0$ 时把速度由 $v$ 变为 $v'$ 的流。则在微扰论各阶：(a) \cref{def:lamet-soft-offlc}的 $S(b_\perp,\mu,Y,Y')$ 等于由色单态 $Q\bar Q$ 插值场在 $-t$ 处产生、在 $t'$ 处湮灭、中间插入 $J$ 的截肢关联函数在 $t,t'\to\infty$ 的极限；(b) 因 $J$ 是等时算符，$S=\lim_{\tau,\tau'\to\infty}W_E/\sqrt{Z_EZ_E}$，其中 $W_E$、$Z_E$ 用欧氏 HQET 计算，其拉氏量为 $\psi_v^\dagger(\ii\tilde v\cdot D_E)\psi_v+\eta_v^\dagger(\ii\tilde v\cdot D_E)\eta_v$，$\tilde v=\gamma(-\ii,-\beta,\mathbf 0_\perp)$，$\ii\tilde v\cdot D_E=\gamma(D_\tau-\ii\beta D_z)$（欧氏内积；常数因子 $\gamma$ 可由场的归一化吸收）；(c) $S$ 只依赖于 $Y+Y'$。（本定理是微扰论中真空关联函数之间的恒等式。）% VF-YZL-2
% Writer note (W26): Ji:2019sxk eq. (13) prints i v~.D_E = D^tau - i beta D^z with v~ = gamma(-i,-beta,0); with the Euclidean
% inner product the left side is gamma(D_tau - i beta D_z) (checked from i v.D, t = -i tau, A_0 = i A_tau); the factor gamma is
% absorbed by the field normalization. Printed here with the factor.
\uses{def:lamet-soft-offlc,lem:yzliu-partial-exp,asm:lamet-aux-renorm}
\end{theorem}

\begin{proofsketch}
（1）静态色源的传播子是 Wilson 线，四个色源的时序关联函数逐图重现两个订书钉形回路，插值场中的横向 Wilson 线给出订书钉的横向连接；（2）$Z(2t)=W(t,t,b_\perp,Y,-Y)$ 是不改变速度的同一关联函数，由\cref{lem:yzliu-partial-exp}，自相互作用与时间演化在 $W/\sqrt{ZZ}$ 中相消（LSZ 型截肢）；（3）静态源与流的乘法可重整性（\cref{asm:lamet-aux-renorm}）使等式对重整化的量成立；（4）初、末态是演化得到的最低态，$J$ 只作用于 $t=0$，把实时间演化换成欧氏演化投影到同一组态，作者在单圈阶显式核对了 (b)\cite[式~(12)--(13)]{Ji:2019sxk}；（5）沿 $z$ 的推动使 $Y\to Y+\eta$、$Y'\to Y'-\eta$。见\cite[式~(1)--(13)]{Ji:2019sxk}，\cite[第~4.1--4.2~节、附录~A]{Liu:2022nnk}，\cite[第~V.C~节]{JiEtAl2021}。
\end{proofsketch}

\paragraph{地位与缺口。}全阶论证 [A]：组合部分依赖\cref{lem:yzliu-partial-exp}（[A]），紫外部分依赖\cref{asm:lamet-aux-renorm}（[A]），欧氏化一步只在单圈阶显式核对；非微扰的形状因子形式是\cref{cor:yzliu-soft-form-factor}（内容等级 [C]）。作者指出的格点困难：静态夸克的倍增问题、流与速度的重整化、大速度色源的噪声\cite{Ji:2019sxk}。此前的欧氏“弯折”软函数\cite{Ji:2014hxa,Ebert:2019okf}只含欧氏直角尖点，单圈时其反常维度与 $\gamma_{HH}$ 一致，但由两圈角依赖尖点反常维度\cite{KorchemskyRadyushkin1987}在欧氏尖点角 $\frac\pi2$（即双曲角 $\theta=\ii\frac\pi2$）处的值，我们得到（以 $(\alpha_s/\pi)^2$ 为单位）% VF-YZL-17
\begin{equation}\label{eq:yzliu-bent-mismatch}
\Gamma^{(2)}_{\mathrm{cusp}}\big(\tfrac\pi2\big)-\gamma^{(2)}_{HH}=C_FC_A\Big(\frac{\zeta_3}{2}-\frac{\pi^2}{24}\Big)\approx0.190\,C_FC_A ,
\end{equation}
$n_f$ 部分相等；这些因子只依赖于 $\mu b_\perp$，故在 $\order{\alpha_s^2}$ 产生任何与 $b_\perp$ 无关的核都无法吸收的 $\ln(\mu^2b_\perp^2)$ 项。这一论证属于文献\cite{Ji:2019sxk}（亦见\cite{Ji:2019ewn}），文献\cite{Ebert:2022fmh}转述了它，\cref{eq:yzliu-bent-mismatch}的数值是本书的计算。据我们所知（INSPIRE 检索，2026 年 10 月），文献\cite{Ji:2019sxk}首次把 Collins 的离光锥软因子表示为与时间无关的矩阵元——快速重夸克对的形状因子；此前的弯折软函数只再现单圈对数。这也纠正了误解 M16：快度发散不是正规化的人为产物——$\ln S$ 随 $Y+Y'$ 线性增长（\cref{prop:yzliu-rapidity-linear}），斜率正是物理形状因子的 Collins--Soper 核。

\begin{example}[单圈离光锥软函数与 $S_r$]\label{ex:yzliu-soft-oneloop}
由\cref{lem:tmd-soft-oneloop}(b)，快度差为 $\mathcal Y$ 的单圈软因子为 $\tilde S^{(1)}=-2a(\mathcal Y\coth\mathcal Y-\sigma)L_b$，$\sigma=1$ 对应含自能的类时线，$\sigma=0$ 对应不含自相互作用的类空线（\cref{def:tmd-soft}）。Ji--Liu--Liu 的软函数对应 $\sigma=1$、$\mathcal Y=Y+Y'$；由 $\coth\mathcal Y=1+2\ee^{-2\mathcal Y}+\cdots$，% VF-TMD-8
\begin{equation}\label{eq:yzliu-soft-oneloop}
S^{(1)}=-2a(\mathcal Y-1)L_b+\order{\mathcal Y\ee^{-2\mathcal Y}}=a\big[2-2(Y+Y')\big]L_b+\order{\mathcal Y\ee^{-2\mathcal Y}} ,
\end{equation}
故 $\tilde K^{(1)}=-2aL_b$，$\mathcal D^{(1)}=2aL_b$，$S_r^{(1)}=-\mathcal D^{(1)}=-\frac{\alpha_sC_F}{\pi}L_b$（\cref{def:lamet-reduced-soft}）。$\sigma=0$ 给出同样的斜率而 $\mathcal D^{(1)}=0$：这正是\cref{def:lamet-soft-offlc}必须采用 Ji--Liu--Liu 归一化（类时线、保留自能、以 $Z$ 因子减除）的原因——单个软函数依赖于正规化方式，只有组合才是物理的。又 $\mu^2\frac{\dd\ln S}{\dd\mu^2}=-\Gamma(\mathcal Y)$，$\Gamma(\mathcal Y)=\frac{\alpha_sC_F}{\pi}(\mathcal Y\coth\mathcal Y-1)\to\mathcal Y\Gamma_{\mathrm{cusp}}+\gamma^{(1)}_{HH}$，$\gamma^{(1)}_{HH}=-\frac{\alpha_sC_F}{\pi}$：于是 $\mu^2\frac{\dd\tilde K}{\dd\mu^2}=-\Gamma_{\mathrm{cusp}}$，$\mu^2\frac{\dd\ln S_r}{\dd\mu^2}=\gamma_{HH}$。
\uses{thm:yzliu-soft-hqet,lem:tmd-soft-oneloop,def:lamet-reduced-soft}
\end{example}

\begin{assumption}[静态色单态 $Q\bar Q$ 对的谱假设 \tagC]\label{asm:yzliu-heavy-pair}
(HP1) 在固定 $\vec b_\perp$ 的静态色单态 $Q\bar Q$ 扇区中存在孤立的基态，剩余能量 $E=\gamma^{-1}V(\vec b_\perp)$，连续谱有能隙（激发态由 Riemann--Lebesgue 引理或 $t\to t(1-\ii0)$ 投影掉）；(HP2) 色源对的插值场以非零、推动不变的交叠 $\Phi(\vec b_\perp)$ 产生基态；(HP3) 规范线自能与色源对的时间演化恰被 $\sqrt{ZZ}$ 消去（LSZ 型截肢；微扰论中这正是部分指数化引理的内容）。地位：[C]——非微扰谱假设，(HP1) 有格点静态势作为证据。文献：\cite{Ji:2019sxk}，\cite[第~4.1--4.2~节、附录~A]{Liu:2022nnk}。被使用于：\cref{cor:yzliu-soft-form-factor}。解除途径：(HP1) 需要静态色单态扇区的能隙（禁闭），(HP2)(HP3) 需要非微扰 HQET 中静态源的 LSZ 型构造。
\uses{def:lamet-soft-offlc}
\end{assumption}

\paragraph{地位与缺口。}猜想级（[C]）：涉及禁闭扇区的谱，没有全阶论证；以它为前提的陈述内容等级都是 [C]，而微扰的\cref{thm:yzliu-soft-hqet}不依赖于它。

\begin{corollary}[软函数是静态色源对的形状因子（非微扰形式） \tagR]\label{cor:yzliu-soft-form-factor}
（$R_{\mathrm{cond}}$：在所列编号假设之下证明完整；结论的地位不高于其中最弱的假设。）前提：\cref{asm:yzliu-heavy-pair}（及 \cref{thm:yzliu-soft-hqet}）。设 $|Q\bar Q,\vec b_\perp\rangle_v$ 为以速度 $v$ 运动、横向间隔 $\vec b_\perp$、剩余动量为零的最低色单态静态色源对，则
\begin{equation}\label{eq:yzliu-soft-ff}
S(b_\perp,Y,Y')={}_{v'}\big\langle Q\bar Q,\vec b_\perp\big|\,J(\vec b_\perp,v,v')\,\big|Q\bar Q,\vec b_\perp\big\rangle_v :
\end{equation}
离光锥软函数是快速色单态静态源对的弹性形状因子。% VF-YZL-2
非微扰地位：\cref{asm:intro-wightman}（MA0）。
\uses{asm:yzliu-heavy-pair,thm:yzliu-soft-hqet}
\end{corollary}

\begin{proof}
由\cref{thm:yzliu-soft-hqet}(a)，$S=\lim W(t,t')/\sqrt{Z(2t)Z(2t')}$，$W=\langle0|O_{v'}(t')\,J(0)\,O^\dagger_v(-t)|0\rangle$，$O_v$ 为色源对的插值场。在 $J$ 前后插入色源对扇区能量本征态 $\{|n\rangle_v\}$（剩余能量 $E_n$）的完备集：$W=\sum_{m,n}\langle0|O_{v'}|m\rangle\ee^{-\ii E'_mt'}\langle m|J|n\rangle\ee^{-\ii E_nt}\langle n|O_v^\dagger|0\rangle$。令 $t,t'\to\infty(1-\ii0)$，由 (HP1) 只有基态存活：$W\to\Phi'^*\Phi\,\ee^{-\ii E'_0t'-\ii E_0t}\,{}_{v'}\langle0|J|0\rangle_v$，$\Phi={}_v\langle0|O_v^\dagger|0\rangle$。$Z(2t)$ 是 $v'=v$ 时的同一关联函数，而 $J(\vec b_\perp,v,v)$ 在色源对态上为恒等，故 $Z(2t)\to|\Phi|^2\ee^{-2\ii E_0t}$，$Z(2t')\to|\Phi'|^2\ee^{-2\ii E'_0t'}$，从而 $W/\sqrt{Z(2t)Z(2t')}\to(\Phi'^*\Phi/|\Phi'\Phi|)\,{}_{v'}\langle0|J|0\rangle_v$。由 (HP2)，$\Phi\ne0$ 且推动不变，$\Phi'=\Phi$，相位因子为 1；由 (HP3)，$\sqrt{ZZ}$ 消去的恰是自能与时间演化，左边的极限即离光锥软函数 $S$ 本身。于是得\cref{eq:yzliu-soft-ff}（参见\cite[式~(8)、(11)]{Ji:2019sxk}）。
\end{proof}

推论的内容等级为 [C]；其格点上的对应物是通过介子形状因子提取本征软函数（\cref{prop:yzliu-meson-ff}）。
